All pointwise parts of the record are finite Boolean combinations of polynomial equalities and strict inequalities with real-algebraic coefficients. A cylindrical sign decomposition is obtained recursively: projection adjoins discriminants, resultants, and leading coefficients; univariate roots are isolated by rational intervals and Sturm sign changes; lifting produces finitely many cells on which every input polynomial has constant sign. The recursion terminates because each projection lowers the number of variables. Evaluating the finitely many signs decides every quantified identity, angle inequality, bi-Lipschitz inequality, and transverse-Jacobian inequality in the proposed statement.

Consider next one certified nonempty link sector. Its conic parametrization and the transverse-Jacobian bounds reduce its Hausdorff mass to
\[
L=\int_V J(v)\,d\mathcal H^h(v),
\qquad 0<c_J\le J(v)\le C_J<\infty,
\]
over a compact semialgebraic cell \(V\). The same sign decomposition triangulates \(V\) into finitely many graph cells. The certificate supplies a common Lipschitz modulus for those graph maps and for \(J\). Subdivision into rational boxes of mesh \(2^{-n}\), interval evaluation of the algebraic graph maps, and the Jacobian bounds therefore give rational numbers \(L_n^-\le L\le L_n^+\) satisfying
\[
0\le L_n^+-L_n^-\le C,2^{-n},
\]
where \(C\) is computed from the finite Lipschitz and Jacobian record. Hence \(L\) is a computable real, uniformly over the finitely many sectors.

If the proposed certificate records rational endpoints \(c_L<C_L\) and the strict inequalities \(c_L<L<C_L\) are true, choose
\[
\epsilon=\tfrac12\min\{L-c_L,C_L-L\}>0.
\]
For all sufficiently large \(n\), the interval \([L_n^-,L_n^+]\) has width below \(\epsilon\), and then it lies inside \((c_L,C_L)\); the procedure accepts. If instead \(L<c_L\) or \(L>C_L\), the same argument eventually produces an interval strictly on the violating side and the procedure rejects. When \(L=c_L\) or \(L=C_L\), interval refinement need not terminate. Thus computability proves the stated two-sided semidecision but not the total decision asserted by the current claim. Active support faces, tangencies, and corners merely add sign cells and their own link sectors, so the argument applies without a transversality assumption.